\chapter{Linear Transformer}

\section{概述}

Linear Transformer 由 Katharopoulos 等人于 2020 年提出，是线性注意力机制的开创性工作之一。该论文首次系统地将注意力计算重写为可分解的核函数形式，实现了关于序列长度的线性复杂度。

\section{从 softmax 到可分解核}

标准 Transformer 的多头自注意力定义为：

\begin{equation}
    \text{Attention}(\mathbf{Q}, \mathbf{K}, \mathbf{V}) = \text{softmax}\left(\frac{\mathbf{Q}\mathbf{K}^\top}{\sqrt{d_k}}\right)\mathbf{V}
\end{equation}

Linear Transformer 的关键洞察是：softmax 操作可以被视为核函数的归一化形式。具体地，softmax 注意力中的相似度函数为：

\begin{equation}
    \text{sim}(\mathbf{q}, \mathbf{k}) = \exp\left(\frac{\mathbf{q}^\top \mathbf{k}}{\sqrt{d_k}}\right)
\end{equation}

但这个核函数无法直接分解为 $\phi(\mathbf{q})^\top \phi(\mathbf{k})$ 的形式。Linear Transformer 使用了一个简单的替代方案：

\begin{equation}
    \phi(\mathbf{x}) = \text{elu}(\mathbf{x}) + 1
\end{equation}

这个特征映射有两个重要性质：

\begin{enumerate}
    \item $\phi(\mathbf{x}) > 0$ 对所有 $\mathbf{x}$ 成立，保证注意力权重非负
    \item 近似捕捉输入之间的正相关性
\end{enumerate}

\section{归一化}

为了替代 softmax 的自动归一化功能，Linear Transformer 引入了显式的归一化项。每个查询位置的注意力输出为：

\begin{equation}
    \mathbf{o}_i = \frac{\sum_{j=1}^N \phi(\mathbf{q}_i)^\top \phi(\mathbf{k}_j) \mathbf{v}_j}{\sum_{j=1}^N \phi(\mathbf{q}_i)^\top \phi(\mathbf{k}_j)}
\end{equation}

利用结合律重写为：

\begin{equation}
    \mathbf{o}_i = \frac{\phi(\mathbf{q}_i)^\top \left(\sum_{j=1}^N \phi(\mathbf{k}_j) \mathbf{v}_j^\top\right)}{\phi(\mathbf{q}_i)^\top \left(\sum_{j=1}^N \phi(\mathbf{k}_j)\right)} = \frac{\phi(\mathbf{q}_i)^\top \mathbf{S}}{\phi(\mathbf{q}_i)^\top \mathbf{z}}
\end{equation}

\section{因果掩码的递推形式}

对于自回归任务，Linear Transformer 利用其递推特性自然实现因果掩码。对于序列位置 $i$：

\begin{equation}
    \mathbf{o}_i = \frac{\phi(\mathbf{q}_i)^\top \mathbf{S}_i}{\phi(\mathbf{q}_i)^\top \mathbf{z}_i}
\end{equation}

其中：

\begin{equation}
    \mathbf{S}_i = \mathbf{S}_{i-1} + \phi(\mathbf{k}_i)\mathbf{v}_i^\top, \quad \mathbf{z}_i = \mathbf{z}_{i-1} + \phi(\mathbf{k}_i)
\end{equation}

初始状态 $\mathbf{S}_0 = \mathbf{0}$，$\mathbf{z}_0 = \mathbf{0}$。

这种递推形式在训练时（非自回归）和推理时（自回归）都有效：

\begin{itemize}
    \item 训练阶段：可以并行计算所有位置的输出，使用累积和操作
    \item 推理阶段：只需维护状态 $(\mathbf{S}_i, \mathbf{z}_i)$，每步 $\mathcal{O}(d^2)$ 计算
\end{itemize}

\section{多头线性注意力}

Linear Transformer 也支持多头机制。对于每个头 $h$：

\begin{equation}
    \text{head}_h = \text{LinearAttention}(\mathbf{Q}_h, \mathbf{K}_h, \mathbf{V}_h)
\end{equation}

其中 $\mathbf{Q}_h, \mathbf{K}_h, \mathbf{V}_h$ 是通过不同的投影矩阵获得的。最终输出为所有头的拼接经过线性投影。

\section{复杂度对比}

\begin{table}[H]
    \centering
    \begin{tabular}{l|cc}
        \toprule
        & 时间复杂度 & 空间复杂度 \\
        \midrule
        标准注意力 & $\mathcal{O}(N^2 d)$ & $\mathcal{O}(N^2 + Nd)$ \\
        Linear Transformer & $\mathcal{O}(Nd^2)$ & $\mathcal{O}(Nd)$ \\
        \bottomrule
    \end{tabular}
    \caption{标准注意力与 Linear Transformer 复杂度对比}
\end{table}

当 $N > d$ 时，Linear Transformer 在时间和空间上都有显著优势。

\section{伪代码}

\begin{algorithm}[H]
    \caption{Linear Transformer 前向传播}
    \begin{algorithmic}[1]
        \Require $\mathbf{Q}, \mathbf{K}, \mathbf{V} \in \mathbb{R}^{N \times d}$, causal: bool
        \Ensure $\mathbf{O} \in \mathbb{R}^{N \times d}$
        \State $\mathbf{Q} \gets \phi(\mathbf{Q})$, $\mathbf{K} \gets \phi(\mathbf{K})$
        \If{causal}
            \State $\mathbf{S} \gets \mathbf{0}_{d \times d}$, $\mathbf{z} \gets \mathbf{0}_d$
            \For{$i = 1$ to $N$}
                \State $\mathbf{S} \gets \mathbf{S} + \mathbf{K}[i]^\top \mathbf{V}[i]$
                \State $\mathbf{z} \gets \mathbf{z} + \mathbf{K}[i]$
                \State $\mathbf{O}[i] \gets (\mathbf{Q}[i]^\top \mathbf{S}) / (\mathbf{Q}[i]^\top \mathbf{z} + \epsilon)$
            \EndFor
        \Else
            \State $\mathbf{S} \gets \mathbf{K}^\top \mathbf{V}$
            \State $\mathbf{z} \gets \sum_{i=1}^N \mathbf{K}[i]$
            \State $\mathbf{O} \gets \mathbf{Q} \mathbf{S} / (\mathbf{Q} \mathbf{z} + \epsilon)$
        \EndIf
        \State \Return $\mathbf{O}$
    \end{algorithmic}
\end{algorithm}
